
\documentclass[10pt]{article}

\usepackage[utf8]{inputenc}
\usepackage[T1]{fontenc}
\usepackage{url}
\usepackage{hyperref}
\usepackage{amsmath}
\usepackage{amsfonts}
\usepackage{amssymb}
\usepackage{graphicx}
\usepackage{float}
\usepackage{lipsum}
\usepackage{multicol}
\usepackage{xcolor}
\usepackage[numbers,square]{natbib}
\usepackage[font=small]{caption}

\usepackage[
    a4paper,
    left=1.50cm,
    right=1.50cm,
    top=1.50cm,
    bottom=1.50cm
]{geometry}

\addtolength{\abovecaptionskip}{-3mm}
\addtolength{\textfloatsep}{-5mm}
\setlength\columnsep{20pt}

\title{Renewable energy expansion in the Baltic States: An input-output analysis of 
the energy trilemma---energy security, equity, and environmental 
sustainability}

\author{
    Daria Mytnik\thanks{McGill University, 
    Montréal, Québec, Canada.}
    \and
    Kakali Mukhopadhyay\thanks{McGill University, 
    Montréal, Québec, Canada.}
    \and
    Carl Swanson\thanks{Environment and Climate Change Canada, 
    Gatineau, Québec, Canada.}
}
\date{}

\begin{document}

\maketitle

\section*{Abstract}

The European Union has accelerated the deployment of renewable energy to reduce greenhouse gas emissions and 
dependence on imported fossil fuels. However, the effects of renewable electricity expansion may differ across 
countries depending on their energy systems and economic structures. Although the Baltic States' energy transitions 
have been widely studied, existing research generally examines the dimensions of the energy trilemma separately. 
In particular, few studies capture how renewable electricity expansion simultaneously affects economy-wide emissions, 
embodied import dependence, and the transmission of electricity costs to households through interconnected 
production sectors. This study addresses this gap by applying an integrated input-output framework to examine 
the energy trilemma implications of renewable electricity expansion in Estonia, Latvia, and Lithuania for 
environmental sustainability, energy security, and household affordability.

We disaggregate the electricity sector in the 2022 OECD input-output table using
GTAP12-Power peak and base-load renewable and fossil electricity generating types and environmentally extend
the table using the OECD's air emission accounts (AEA). We then construct
renewable electricity transition scenarios aligned with the National Energy and Climate Plans of 
the Baltic States. Disaggregating the electricity sector into all 11 GTAP12-Power generating types and
keeping the split between peak and base-load generation enables 
the estimation of economy-wide effects on emissions, import dependence, 
household affordability, capacity adequacy, and macroeconomic indicators under a 2030 business-as-usual (BAU) baseline and 
alternative renewable energy scenarios. We assess capacity adequacy through 
country-specific sensitivities benchmarked against published reliability criteria and 
transmission system operator projections, applying Estonia's 10\% capacity-margin criterion and 
considering Latvia's 3-hour and Lithuania's 8-hour annual loss-of-load expectation (LOLE) standards.

Our results indicate that renewable electricity expansion generates limited aggregate macroeconomic effects, 
with GDP changes remaining below 0.05\% and employment increasing by up to 0.12\% under the RES-E scenario. 
Environmental benefits are greatest in Estonia, where economy-wide CO$_2$-equivalent emissions decline by 15.6\%, 
compared with approximately 0.7\% in Latvia and Lithuania, primarily reflecting the displacement of carbon-intensive oil shale generation.
Energy security also improves, with embodied import dependence in domestic electricity supply declining by 47.2\% in Lithuania, 
15.3\% in Estonia, and 2.9\% in Latvia. Affordability impacts remain relatively modest: under the combined electricity-price-pressure 
scenario, household price indices increase by 0.78\% in Latvia, 1.32\% in Estonia, and 1.44\% in Lithuania. 
Overall, our findings highlight the contribution of renewable electricity expansion to the energy trilemma 
and to progress toward the UN Sustainable Development Goals, particularly SDG 7 (Affordable and Clean Energy), 
SDG 9 (Industry, Innovation and Infrastructure), and SDG 13 (Climate Action). However, uneven 
effects across countries and energy-trilemma dimensions underscore the need for complementary policies 
to advance the Paris Agreement's climate objectives while safeguarding household affordability.
% \bibliography{bibliography}
% \bibliographystyle{plainnat}

\end{document}
